\section{Результаты моделирования}

\subsection{Параметры модели}

Моделирование проведено для радиального подшипника скольжения геометрии~\varGeomLetter{} ($R = \varRmm$~мм, $c = \varCmkm$~мкм, $L = \varLmm$~мм) с \varDimpleTypeGenitive{} микрорельефом. Все исходные данные сведены в таблицу~\ref{tab:input_data}.

\begin{table}[H]
\small
\caption{Исходные данные для расчёта}
\label{tab:input_data}
\begin{tabularx}{\textwidth}{|X|c|c|}
\hline
\centering\textbf{Параметр} & \textbf{Значение} & \textbf{Размерность} \\
\hline
Радиус вала $R$ & \varRmm & мм \\
\hline
Радиальный зазор $c$ & \varCmkm & мкм \\
\hline
Длина подшипника $L$ & \varLmm & мм \\
\hline
Относительная длина $L/D$ & \varLD & -- \\
\hline
Полуось по оси $Z$, $a$ & \varDimpleAmm & мм \\
\hline
Полуось по оси $\varphi$, $b$ & \varDimpleBmm & мм \\
\hline
Глубина углубления $h_p$ & \varDimpleHpmkm & мкм \\
\hline
Относительная глубина $h_p/c$ & \varDimpleHpRatio & -- \\
\hline
Число углублений по окружности $N_\varphi$ & \varNphi & шт. \\
\hline
Число рядов по длине $N_Z$ & \varNZ & шт. \\
\hline
Область текстурирования & \varPhiRange & -- \\
\hline
Частота вращения $n$ & \varN & об/мин \\
\hline
Вязкость смазки $\eta$ & \varEta & Па$\cdot$с \\
\hline
Диапазон эксцентриситета $\varepsilon$ & 0{,}05-0{,}80 & -- \\
\hline
Размер сетки & $500 \times 500$ & -- \\
\hline
Параметр релаксации $\omega$ & 1{,}5 & -- \\
\hline
Критерий сходимости $\delta$ & $< 10^{-5}$ & -- \\
\hline
Макс. число итераций & 50000 & -- \\
\hline
\end{tabularx}
\end{table}

При заданной частоте вращения линейная скорость вала равна $U = \pi D n / 60 = \pi \cdot 0{,}070 \cdot 2980 / 60 \approx 10{,}93$~м/с. Моделирование выполнено для 15 значений эксцентриситета в указанном диапазоне.

\subsection{Распределение зазора и давления}

На рисунке~\ref{fig:pressure_2d} приведено распределение безразмерного давления $P(\varphi)$ в среднем сечении ($Z = 0$) при $\varepsilon = 0{,}6$ для гладкого и текстурированного подшипников.

\begin{figure}[H]
\centering
\includegraphics[width=\textwidth]{figures/pressure_2d.png}
\caption{Распределение давления $P(\varphi)$ при $Z=0$ для гладкого и текстурированного подшипников ($\varepsilon = 0{,}6$)}
\label{fig:pressure_2d}
\end{figure}

Кривая давления гладкого подшипника имеет характерную колоколообразную форму с максимумом вблизи области минимального зазора ($\varphi \approx \pi$). В расширяющейся части зазора давление обращается в нуль вследствие кавитации, а область положительного давления охватывает приблизительно половину окружности. Введение цилиндрического микрорельефа приводит к появлению выраженных локальных возмущений на кривой давления. Характерной особенностью цилиндрического профиля является наличие резких перепадов давления на границах углублений, обусловленных скачкообразным изменением зазора. В отличие от эллипсоидального профиля, где глубина меняется плавно, цилиндрическое углубление создаёт ступенчатый переход, что формирует более интенсивные пики давления на входных кромках. Несмотря на это, общий уровень давления текстурированного подшипника существенно выше, обеспечивая значительный прирост несущей способности.

На рисунке~\ref{fig:clearance_2d} изображён профиль зазора $H(\varphi)$ в среднем сечении при $\varepsilon = 0{,}6$.

\begin{figure}[H]
\centering
\includegraphics[width=\textwidth]{figures/clearance_2d.png}
\caption{Распределение зазора $H(\varphi)$ при $Z=0$ для подшипника с \varDimpleTypeGenitive{} углублениями ($\varepsilon = 0{,}6$)}
\label{fig:clearance_2d}
\end{figure}

Профиль зазора следует закону $H = 1 + \varepsilon\cos\varphi$, на фоне которого видны ступенчатые увеличения, характерные для цилиндрических углублений. Глубина каждого углубления в безразмерных единицах составляет $\Delta H_{\max} = h_p/c = 0{,}2$. Переходы на кромках углублений имеют вертикальный характер, что отличает цилиндрический профиль от плавных эллипсоидальных форм.

Трёхмерные распределения зазора и давления при $\varepsilon = 0{,}6$ показаны на рисунках~\ref{fig:clearance_3d} и~\ref{fig:pressure_3d}.

\begin{figure}[H]
\centering
\includegraphics[width=\textwidth]{figures/clearance_3d.png}
\caption{Трёхмерное распределение зазора $H(\varphi, Z)$ ($\varepsilon = 0{,}6$)}
\label{fig:clearance_3d}
\end{figure}

Рисунок~\ref{fig:clearance_3d} демонстрирует регулярный массив цилиндрических углублений: 8 по окружности и 11 рядов по длине подшипника, размещённых в зоне $90^\circ - 270^\circ$.

\begin{figure}[H]
\centering
\includegraphics[width=\textwidth]{figures/pressure_3d.png}
\caption{Трёхмерное распределение давления $P(\varphi, Z)$ ($\varepsilon = 0{,}6$)}
\label{fig:pressure_3d}
\end{figure}

Из рисунка~\ref{fig:pressure_3d} видно, что каждое цилиндрическое углубление формирует резкий подъём давления на входной кромке, где поток сталкивается со ступенчатым расширением зазора. Этот микрогидродинамический эффект при цилиндрическом профиле проявляется более отчётливо, чем при плавных формах, поскольку резкие границы углубления создают выраженные локальные градиенты давления. На выходной кромке наблюдается спад давления, однако кавитационное ограничение ($P \geq 0$) обеспечивает положительный суммарный вклад каждого углубления в несущую способность.

При рассмотрении трёхмерных распределений следует выделить ряд характерных закономерностей. Максимальное давление формируется в среднем сечении ($Z = 0$), а к торцам подшипника ($Z = \pm 1$) давление спадает до нуля из-за торцевых утечек. Микрорельеф оказывает наибольшее влияние в центральной области, где фоновое давление достигает максимума. Вблизи торцов эффект текстурирования ослабевает вследствие низкого уровня окружающего давления.

\subsection{Сравнительный анализ интегральных показателей}

В таблицах~\ref{tab:smooth_results} и~\ref{tab:textured_results} приведены интегральные характеристики гладкого подшипника и подшипника с \varDimpleTypeGenitive{} углублениями для характерных значений эксцентриситета.

\begin{table}[H]
\small
\centering
\caption{Характеристики гладкого подшипника}
\label{tab:smooth_results}
\begin{tabular}{|c|c|c|c|}
\hline
$\boldsymbol{\varepsilon}$ & $\boldsymbol{F}$\textbf{, Н} & $\boldsymbol{\mu}$ & $\boldsymbol{Q}$\textbf{, л/с} \\
\hline
0{,}05 & 259{,}5 & 0{,}2291 & 0{,}2398 \\
\hline
0{,}16 & 835{,}7 & 0{,}0722 & 0{,}2395 \\
\hline
0{,}26 & 1478{,}5 & 0{,}0420 & 0{,}2389 \\
\hline
0{,}37 & 2261{,}3 & 0{,}0287 & 0{,}2380 \\
\hline
0{,}48 & 3286{,}8 & 0{,}0211 & 0{,}2370 \\
\hline
0{,}59 & 4764{,}4 & 0{,}0160 & 0{,}2360 \\
\hline
0{,}69 & 7222{,}5 & 0{,}0120 & 0{,}2351 \\
\hline
0{,}80 & 12218{,}3 & 0{,}0087 & 0{,}2344 \\
\hline
\end{tabular}
\end{table}

\begin{table}[H]
\small
\centering
\caption{Характеристики подшипника с \varDimpleTypeGenitive{} углублениями}
\label{tab:textured_results}
\begin{tabular}{|c|c|c|c|}
\hline
$\boldsymbol{\varepsilon}$ & $\boldsymbol{F}$\textbf{, Н} & $\boldsymbol{\mu}$ & $\boldsymbol{Q}$\textbf{, л/с} \\
\hline
0{,}05 & 635{,}1 & 0{,}0922 & 0{,}2448 \\
\hline
0{,}16 & 1196{,}4 & 0{,}0495 & 0{,}2444 \\
\hline
0{,}26 & 1925{,}0 & 0{,}0317 & 0{,}2437 \\
\hline
0{,}37 & 2868{,}3 & 0{,}0222 & 0{,}2428 \\
\hline
0{,}48 & 4251{,}6 & 0{,}0160 & 0{,}2418 \\
\hline
0{,}59 & 6583{,}4 & 0{,}0113 & 0{,}2407 \\
\hline
0{,}69 & 11407{,}7 & 0{,}0074 & 0{,}2396 \\
\hline
0{,}80 & 25685{,}4 & 0{,}0040 & 0{,}2384 \\
\hline
\end{tabular}
\end{table}

С увеличением эксцентриситета нагрузочная способность $F$ растёт нелинейно для обоих вариантов подшипника. Подшипник с цилиндрическим микрорельефом демонстрирует значительно более высокую нагрузочную способность во всём диапазоне $\varepsilon$. Особенно существенно различие при больших эксцентриситетах: при $\varepsilon = 0{,}80$ нагрузочная способность текстурированного подшипника (25685{,}4~Н) превосходит гладкий аналог (12218{,}3~Н) более чем вдвое. Столь значительный эффект объясняется тем, что цилиндрические углубления с резкими кромками создают интенсивные локальные перепады давления, генерирующие выраженный микрогидродинамический подъём.

Коэффициент трения $\mu$ уменьшается с ростом $\varepsilon$ для обоих вариантов, причём текстурированный подшипник имеет более низкие значения $\mu$ во всём диапазоне. Расход смазки $Q$ практически не зависит от эксцентриситета и несколько выше у текстурированного подшипника вследствие увеличения эффективного объёма смазочного слоя в области углублений.

Количественное сопоставление характеристик при $\varepsilon = 0{,}6$ выполнено в таблице~\ref{tab:comparison}.

\begin{table}[H]
\small
\centering
\caption{Сравнение характеристик при $\varepsilon = 0{,}6$}
\label{tab:comparison}
\begin{tabular}{|l|c|c|c|}
\hline
\textbf{Характеристика} & \textbf{Гладкий} & \textbf{С углублениями} & $\boldsymbol{\Delta}$\textbf{, \%} \\
\hline
Нагрузка $F$, Н & 5028{,}5 & 7001{,}5 & $+39{,}24$ \\
\hline
Коэфф. трения $\mu$ & 0{,}0154 & 0{,}0108 & $-30{,}03$ \\
\hline
Расход $Q$, л/с & 0{,}2359 & 0{,}2405 & $+1{,}97$ \\
\hline
\end{tabular}
\end{table}

Применение цилиндрических углублений обеспечивает прирост несущей способности подшипника на 39{,}24\%, что обусловлено микрогидродинамическим эффектом на резких кромках углублений. Коэффициент трения снижается на 30{,}03\% благодаря возрастанию средней толщины смазочного слоя в зонах текстурирования. Увеличение расхода смазки на 1{,}97\% содействует дополнительному теплоотводу и не предъявляет чрезмерных требований к системе маслоснабжения.

На рисунках~\ref{fig:load_vs_eps}-\ref{fig:flow_vs_eps} показаны зависимости интегральных характеристик от эксцентриситета для обоих вариантов подшипника.

\begin{figure}[H]
\centering
\includegraphics[width=\textwidth]{figures/load_vs_epsilon.png}
\caption{Нагрузочная способность $F(\varepsilon)$ для гладкого и текстурированного подшипников}
\label{fig:load_vs_eps}
\end{figure}

Зависимость $F(\varepsilon)$ текстурированного подшипника расположена значительно выше кривой гладкого подшипника на всём интервале эксцентриситетов. Разрыв между кривыми нарастает с увеличением $\varepsilon$: при $\varepsilon = 0{,}05$ разница составляет около 376~Н, тогда как при $\varepsilon = 0{,}80$ она превышает 13400~Н. Нелинейный рост обусловлен кубической зависимостью давления от обратной толщины зазора ($P \sim h^{-3}$), вследствие чего при малых зазорах эффект каждого углубления многократно усиливается.

\begin{figure}[H]
\centering
\includegraphics[width=\textwidth]{figures/friction_vs_epsilon.png}
\caption{Коэффициент трения $\mu(\varepsilon)$ для гладкого и текстурированного подшипников}
\label{fig:friction_vs_eps}
\end{figure}

Для текстурированного подшипника коэффициент трения ниже гладкого при всех значениях $\varepsilon$. Кривые монотонно убывают и асимптотически стремятся к нулю при $\varepsilon \to 1$. Максимальный абсолютный выигрыш приходится на малые эксцентриситеты: при $\varepsilon = 0{,}05$ разность $\Delta\mu = 0{,}137$, тогда как при $\varepsilon = 0{,}80$ она составляет лишь $\Delta\mu = 0{,}005$. Относительное снижение трения остаётся значительным на протяжении всего диапазона эксцентриситетов.

\begin{figure}[H]
\centering
\includegraphics[width=\textwidth]{figures/flow_vs_epsilon.png}
\caption{Расход смазки $Q(\varepsilon)$ для гладкого и текстурированного подшипников}
\label{fig:flow_vs_eps}
\end{figure}

Расход смазки текстурированного подшипника несколько превышает расход гладкого. Разница находится в пределах 1{,}5-2\% и мало изменяется при варьировании эксцентриситета.

Проведённый анализ свидетельствует о том, что цилиндрическое текстурирование позволяет одновременно существенно повысить нагрузочную способность и уменьшить коэффициент трения при умеренном приросте расхода смазки. Эффект проявляется тем сильнее, чем выше эксцентриситет, что делает данный тип микрорельефа перспективным для тяжелонагруженных подшипниковых узлов.

\subsection*{Выводы по главе 3}
\phantomsection
\addcontentsline{toc}{subsection}{Выводы по главе 3}

Проведён расчёт рабочих характеристик радиального подшипника скольжения с \varDimpleTypeGenitive{} микрорельефом для 15 значений эксцентриситета. Построены поля давления и зазора, определены зависимости нагрузочной способности, коэффициента трения и расхода смазки от эксцентриситета. Сравнение при $\varepsilon = 0{,}6$ выявило прирост нагрузочной способности на $39{,}24$\%, снижение коэффициента трения на $30{,}03$\% и увеличение расхода на $1{,}97$\%, что свидетельствует о высокой эффективности цилиндрического текстурирования.
